\documentclass[spanish,openany]{report}
% \usepackage{darkmode} \enabledarkmode
\usepackage{wrapfig}
\usepackage{enumitem}
\input{~/git/Clase/template/preamble.tex}
\input{~/git/Clase/template/macros.tex}
\input{~/git/Clase/template/letterfonts.tex}

\date{\today}
% \hbadness=99999
\usepackage{microtype}
% \addbibresource{./bibliografia.bib}
\usepackage{titling}
\usepackage{afterpage}


% Header and footer
% \fancypagestyle{plain}{
%     \fancyhf{} %Clear Everything.
%     \fancyfoot[C]{Página \thepage\ de \pageref{LastPage} }%Page Number
%     \renewcommand{\headrulewidth}{1pt} %0pt for no rule, 2pt thicker etc...
%     \renewcommand{\footrulewidth}{1pt}
%     \fancyfoot[R]{\small{Daniel Carbajales Soto}}
%     \fancyhead[R]{\small{Cinta Transportadora con 3 Velocidades}}
%     \fancyhead[L]{\small{\today}}
% }
% \pagestyle{plain}

\begin{document}

%----------------------------------------------------------------------------------------
%	                            PORTADA
%----------------------------------------------------------------------------------------
% \makeportada{}{Daniel Carbajales Soto - 1° CFGSARI}{}

%----------------------------------------------------------------------------------------
%                                    HEADER 
%----------------------------------------------------------------------------------------
% \header{}{Daniel Carbajales Soto}{1° CFGSARI}
%----------------------------------------------------------------------------------------
%                               TABLE OF CONTENTS
%----------------------------------------------------------------------------------------
% \maketoc
%----------------------------------------------------------------------------------------
%	                            TABLES
%----------------------------------------------------------------------------------------
% this is just a demo in case i need to make a table
%
%
%
% \begin{tabularx}{\linewidth}{|C{2cm}|Y|Y|C{2cm}|}
% \hline
% \makecell{Header 1} & \makecell{Header 2 \\ with linebreak} & \makecell{Header 3} & \makecell{Header 4} \\
% \hline
% Row 1 & This is a long cell text that automatically wraps in the Y column & Another long cell & Short \\
% \hline
% Row 2 & \multirow{2}{*}{Multirow text} & More text & Short \\
% \cline{1-1} \cline{3-4}
% Row 3 & & Extra text & Short \\
% \hline
% \end{tabularx}
%----------------------------------------------------------------------------------------
% 					DOCUMENT
%----------------------------------------------------------------------------------------
\subsection*{Ejercicio 6}
Dos corrientes alternas senoidales están desfasadas 20º. Sabiendo que la frecuencia es de 50 Hz. Calcular:


a) El período\hspace{5em}  b) El tiempo de desfase.

\incfig[0.25]{desfase-tiempo}

\boldit{a)} El período es de: \[t = \frac{1}{50\si{\hertz}} \to  0'02\si{\second}\].

\boldit{b)} El tiempo de desfase es de:

\[t_{\text{desfase}} = \frac{t}{\upvarphi_{360}}\times\upvarphi_{20} \to 
\frac{\frac{1}{50}}{\ang{360}}\times\ang{20} = 0'0011\si{\second}
\]

%%%%%%%%%%%%%%%%%%%%%%%%%%%%%%%%%%%%%%%%%%%%%%%%%%%%%%%%%%
%%%%%%%%%%%%%%%%%%%% bibliography %%%%%%%%%%%%%%%%%%%%%%%%
% \nocite{*}                                 %%%%%%%%%%%%%
% \printbibliography                         %%%%%%%%%%%%%
%%%%%%%%%%%%%%%%%%%%%%%%%%%%%%%%%%%%%%%%%%%%%%%%%%%%%%%%%%
%%%%%% Last page in case there is no bibliography %%%%%%%%
% \label{Last Page}                           %%%%%%%%%%%%%
%%%%%%%%%%%%%%%%%%%%%%%%%%%%%%%%%%%%%%%%%%%%%%%%%%%%%%%%%%

\end{document}
